\documentclass[titlepage,a4paper]{jsarticle}
\usepackage{../../../sty/import}% 各種パッケージインポート
\usepackage{../../../sty/title}% タイトルページの変更

%% タイトルページの変数
% レポートタイトル
\title{量子・原子力統合特論第4回目レポート}
% 提出日
\expdate{\today}
% 科目名
\subject{量子・原子力統合特論}
% 分野
\class{情報経営システム工学分野}
% 学年
\grade{M1}
% 学籍番号
\mynumber{24336488}
% 記述者
\author{本間三暉}

\begin{document}
% titleページ作成
\maketitle

\section*{高速炉の重要性}

原子力エネルギーを持続可能なエネルギー源として利用するためには，核燃料サイクルの中で高速炉を利用する意義が大きい．軽水炉と高速炉の違いは，単に炉型の違いではなく，中性子のエネルギーと核燃料資源の利用効率の違いにある．

まず，天然ウランの同位体組成を見ると，核分裂しやすい
\(^{235}\mathrm{U}\)は0.7\%しか含まれず，大部分の99.3\%はそのままでは核分裂しにくい
\(^{238}\mathrm{U}\)である．軽水炉では，この天然ウランを濃縮して
\(^{235}\mathrm{U}\)を3〜5\%程度に高めた濃縮ウランを燃料として用いる．一方，濃縮後に残る劣化ウランは
\(^{235}\mathrm{U}\)が0.2\%，\(^{238}\mathrm{U}\)が99.8\%であり，軽水炉だけでは天然ウランの大部分を占める
\(^{238}\mathrm{U}\)を十分に利用できない．これに対し，高速炉では
\(^{238}\mathrm{U}\)を中性子吸収によって
\(^{239}\mathrm{Pu}\)へ転換し，核分裂性物質として再利用できる．また，MOX燃料ではPuを5〜20\%程度含む燃料を利用できるため，UとPuを循環させる核燃料サイクルと相性がよい．

この違いは，中性子の使い方から説明できる．核分裂で放出された中性子のうち，1個は次の核分裂を起こして連鎖反応を維持するために必要である．さらに，一部の中性子は燃料親物質以外に吸収されたり，炉心外へ漏れたりする．資料では，吸収された中性子1個あたりの放出中性子数を
\(\eta\)とし，燃料親物質の転換に使える中性子数を

\[
\eta - 1 - L
\]

として表している．ここで，1は連鎖反応維持に使われる中性子，\(L\)は親物質以外への吸収および系外への漏れを表す．
したがって，\(\eta\)が大きいほど，連鎖反応を維持した後に\(^{238}\mathrm{U}\)をPuへ転換するための余剰中性子が増える．

軽水炉では水によって中性子を減速し，熱中性子を利用して核分裂を起こす．
これは\(^{235}\mathrm{U}\)の核分裂を起こしやすくする点では有効である．
しかし，熱中性子炉では，放出された中性子の多くが連鎖反応の維持や吸収・漏れで失われるため，\(^{238}\mathrm{U}\)を十分にPuへ転換する余裕が小さい．
一方，高速炉では中性子を減速せず，高速中性子を利用する．
高速中性子領域ではPu核種などの\(\eta\)が大きくなるため，連鎖反応を維持した後も，親物質である\(^{238}\mathrm{U}\)の転換に使える中性子を確保しやすい．
このため，高速炉では転換比を1に近づけるだけでなく，1を超える増殖比の領域を目指すことができる．

資源利用効率を定量的に比較すると，高速炉の重要性はさらに明確である．
資料では，天然ウラン利用効率が，軽水炉ワンススルーで約0.5\%，軽水炉プルサーマルで約0.75\%，高速増殖炉で約60\%程度と示されている．
したがって，高速増殖炉は軽水炉ワンススルーと比較して，

\[
\frac{60}{0.5}=120
\]

より，約120倍のウラン利用効率を持つ．また，プルサーマルと比較しても，

\[
\frac{60}{0.75}=80
\]

より，約80倍の利用効率となる．
これは，軽水炉では主に0.7\%しか存在しない\(^{235}\mathrm{U}\)に依存するのに対し，高速炉では99.3\%を占める\(^{238}\mathrm{U}\)をPuへ転換して利用できるためである．

以上より，軽水炉は熱中性子を利用して
\(^{235}\mathrm{U}\)を効率よく核分裂させる炉であるが，天然ウラン全体の利用率は低い．一方，高速炉は高速中性子を利用し，
\(^{238}\mathrm{U}\)を\(^{239}\mathrm{Pu}\)へ転換することで，燃料を消費するだけでなく新たな核分裂性物質を生み出すことができる．
天然ウラン利用効率が軽水炉ワンススルーの約0.5\%から高速増殖炉の約60\%へ向上することを考えると，高速炉はウラン資源の利用可能期間を大きく延ばす技術である．
したがって，高速炉は資源論の立場から，原子力エネルギーを長期的かつ持続可能に利用するために極めて重要な役割を持つ．

% 参考文献
\begin{thebibliography}{99}
\bibitem{a}量子原子力統合工学概論 講義スライド
\end{thebibliography}

\end{document}